\appendix
\section{Formal Proofs for Averaging as Spectral Approximation}
\label{app:proofs}

This appendix provides rigorous proofs supporting
Section~\ref{sec:averaging_spectral}.
To avoid circularity, we first prove results with respect to a fixed linguistic
subspace $\mathcal L$, and then relate it to the empirical SVD spike subspace.

\subsection{Setup}

Assume token representations satisfy
\[
x_i = \ell + s_i,
\qquad
X = \mathbf{1}\ell^\top + S,
\]
where $\ell \in \mathcal L \subset \mathbb{R}^H$ and $s_i$ are semantic residuals.

Let $P_{\mathcal L}$ and $P_{\mathcal L^\perp}$ denote orthogonal projectors onto
$\mathcal L$ and its complement.

\subsection{Assumptions}

\begin{assumption}[Low-rank linguistic structure]
$\ell \in \mathcal L$, where $\dim(\mathcal L)=d_\ell \ll H$.
\end{assumption}

\begin{assumption}[Semantic second-moment bound]
There exists $\sigma^2$ such that
\[
\mathbb{E}\|s_i\|_2^2 \le \sigma^2.
\]
\end{assumption}

\begin{assumption}[Semantic incoherence across sequences]
For tokens drawn from different sequences (or sufficiently distant positions),
the normalized semantic inner product is small. Specifically, for $i \neq j$,
\begin{equation}
\label{eq:assump_cos_incoh}
\left|
\mathbb{E}\!\left[
\frac{s_i^\top s_j}{\|s_i\|_2\,\|s_j\|_2}
\right]
\right|
\le \varepsilon_{\mathrm{inc}},
\end{equation}
for a small constant $\varepsilon_{\mathrm{inc}} \ll 1$.
Moreover, in the tail eigenbasis used in Section~2.4, the normalized coefficient
vectors are also weakly aligned across sequences, i.e.,
$\left|\mathbb{E}[\cos(c_i, c_j)]\right| \le \varepsilon_{\mathrm{inc}}$, where
$c_i$ denotes the tail-basis coefficient vector of $s_i$.
\end{assumption}

This assumption directly matches the empirical cosine similarity measurements in
Section~2.4.

\subsection{Mean Removal Identity}

\begin{lemma}
Let $\bar{s}_N = \frac{1}{N}\sum_{i=1}^N s_i$.
Then
\[
\mu = \ell + \bar{s}_N,
\qquad
X_{\mathrm{avg}} = S - \mathbf{1}\bar{s}_N^\top.
\]
\end{lemma}

\subsection{Concentration of Semantic Mean}

\begin{lemma}
With probability at least $1-\delta$,
\[
\|\bar{s}_N\|_2 \le \frac{\sigma}{\sqrt{N\delta}}.
\]
\end{lemma}

\subsection{Tail Preservation under Mean Removal}
\label{app:tail_preservation}

\begin{theorem}[Tail perturbation is rank-one and controlled by mean leakage]
\label{thm:app_tail_additive}
For any projector $P$ on $\mathbb{R}^H$ and any batch mean $\mu$,
\begin{equation}
\label{eq:app_tail_exact}
X_{\mathrm{avg}}(I-P) - X(I-P) = -\mathbf{1}\mu^\top(I-P),
\end{equation}
and therefore
\begin{equation}
\label{eq:app_tail_bound}
\|X_{\mathrm{avg}}(I-P) - X(I-P)\|_F
=
\|\mathbf{1}\mu^\top(I-P)\|_F
=
\sqrt{N}\,\|(I-P)\mu\|_2.
\end{equation}
In particular, when $P=P_k$ is the empirical spike projector and $\mu$ is well
captured by $\mathrm{range}(P_k)$, tail perturbation is small.
\end{theorem}

\begin{proof}
Equation~\eqref{eq:app_tail_exact} follows directly from
$X_{\mathrm{avg}} = X - \mathbf{1}\mu^\top$.
For \eqref{eq:app_tail_bound}, use the Frobenius norm identity
$\|\mathbf{1}w^\top\|_F = \|\mathbf{1}\|_2\|w\|_2 = \sqrt{N}\|w\|_2$ with
$w=(I-P)\mu$.
\end{proof}

\begin{corollary}[Relative tail preservation under a tail-energy lower bound]
\label{cor:app_tail_ratio}
Assume $\|X(I-P)\|_F^2 \ge cN$ for some constant $c>0$.
Then
\[
\frac{\|X_{\mathrm{avg}}(I-P)\|_F^2}{\|X(I-P)\|_F^2}
=
1 + O\!\left(\frac{\|(I-P)\mu\|_2}{\sqrt{N}}\right).
\]
\end{corollary}

\begin{proof}
Let $A=X(I-P)$ and $E=\mathbf{1}\mu^\top(I-P)$.
Then $X_{\mathrm{avg}}(I-P)=A-E$ and
\[
\bigl|\|A-E\|_F^2 - \|A\|_F^2\bigr|
\le 2\|A\|_F\|E\|_F + \|E\|_F^2.
\]
Divide both sides by $\|A\|_F^2$ and use $\|E\|_F=\sqrt{N}\|(I-P)\mu\|_2$ and
$\|A\|_F^2\ge cN$.
\end{proof}

\paragraph{Definition (Semantic leakage into the linguistic subspace).}
We say semantic leakage into $\mathcal L$ is small if
\[
\mathbb{E}\|P_{\mathcal L} s_i\|_2^2 \le \rho^2,
\qquad \rho \ll \|\ell\|_2.
\]

\subsection{Spike Suppression}

\begin{theorem}
Mean removal eliminates the coherent linguistic component:
\[
XP_{\mathcal L} = \mathbf{1}\ell^\top + SP_{\mathcal L},
\qquad
X_{\mathrm{avg}}P_{\mathcal L} = SP_{\mathcal L} - \mathbf{1}\bar{s}_N^\top P_{\mathcal L}.
\]
If semantic leakage into $\mathcal L$ is small in the above sense, then
\[
\|X_{\mathrm{avg}}P_{\mathcal L}\|_F^2 \ll \|XP_{\mathcal L}\|_F^2.
\]
\end{theorem}

\subsection{Relation to the Empirical SVD Spike Subspace}

Let $P_k$ denote the projector onto the top-$k$ right singular vectors of $X$.
Empirically (Section~2.3),
\[
\frac{\|P_k \mu\|_2^2}{\|\mu\|_2^2} \approx 1
\]
for modest $k$, indicating that $\mathcal L \subseteq \mathrm{range}(P_k)$ up to
small error.
Under this condition, the fixed-subspace guarantees above apply to the empirical
spike subspace used in practice.

\paragraph{Remark.}
Because averaging removes exactly one rank-one component, its relative spike
reduction decreases as $k$ grows.
This behavior matches the empirical spike-reduction curves in Section~3 and is a
direct consequence of the conservative nature of averaging.

\section{Gradient Decomposition Under Mean-Centered Activations}
\label{app:grad_decomposition}

This appendix clarifies how mean-centering activations interacts with gradient
computation in transformer layers, and why removing a shared mean component does
not eliminate token-specific learning signal, even though per-token gradients
are outer products.
The discussion here is explanatory and complements the formal analysis in
Section~4.

\paragraph{Per-token gradient structure.}
Consider a linear projection in a transformer layer,
\[
y = Wx, \qquad W \in \mathbb{R}^{d \times d}, \; x \in \mathbb{R}^d,
\]
and let $\ell_t$ denote the loss associated with token $t$.
The gradient of $\ell_t$ with respect to $W$ is
\[
g_t \triangleq \nabla_W \ell_t = \delta_t x_t^\top,
\]
where $\delta_t = \nabla_{y_t} \ell_t$ is the backpropagated error signal.
The batch gradient used by SGD is the average of these outer products,
\[
G = \frac{1}{N}\sum_{t=1}^N g_t = \mathbb{E}[\delta x^\top].
\]

\paragraph{Decomposition into mean and residual components.}
Decompose the activation into its mean and residual components,
\[
x = \mu_x + \tilde{x},
\qquad
\mu_x = \mathbb{E}[x],
\qquad
\mathbb{E}[\tilde{x}] = 0.
\]
Substituting into the batch gradient yields the exact identity
\begin{equation}
\label{eq:grad_decomp}
\mathbb{E}[\delta x^\top]
=
\underbrace{\mathbb{E}[\delta]\mu_x^\top}_{\text{rank-1, shared}}
+
\underbrace{\mathbb{E}[\delta \tilde{x}^\top]}_{\text{cross-covariance}}.
\end{equation}

\noindent
Equation~\eqref{eq:grad_decomp} separates the batch gradient into two conceptually
distinct components:
\begin{itemize}
  \item the mean term $\mathbb{E}[\delta]\mu_x^\top$, which is a shared, low-rank
  update arising from components that are coherent across tokens (the
  grammar-driven spike); and
  \item the cross-covariance term $\mathbb{E}[\delta \tilde{x}^\top]$, which
  captures correlations between error signals and token-wise variations and
  constitutes the learning-relevant semantic signal.
\end{itemize}

\paragraph{Effect of mean-centering.}
If activations are mean-centered, $x' = x - \mu_x = \tilde{x}$, then the batch
gradient becomes
\[
\mathbb{E}[\delta (x')^\top] = \mathbb{E}[\delta \tilde{x}^\top].
\]
Thus, mean-centering removes only the rank-one mean--mean contribution
$\mathbb{E}[\delta]\mu_x^\top$ and leaves the cross-covariance term intact.
This explains why averaging suppresses a coherent, spike-like component of the
gradient while preserving token-specific learning signal.

\paragraph{Relation to second-order structure and adaptive optimizers.}
Although SGD updates use the first moment of per-token gradients, the learning
signal associated with semantic variation is expressed through correlations
between $\delta_t$ and $\tilde{x}_t$ across tokens.
Moreover, adaptive optimizers such as Adam and AdamW explicitly maintain
second-moment estimates of gradients,
\[
v \approx \mathbb{E}[g_t \odot g_t],
\]
which decompose into contributions from shared components and from token-wise
variability.
A large shared component inflates these second moments and suppresses effective
learning rates in semantically meaningful directions.
Mean-centering reduces this shared contribution while leaving the variability
term largely unchanged, which is consistent with the empirical improvements
observed in optimization and low-bit quantization.

\paragraph{Summary.}
Although per-token gradients are outer products, mean-centering activations
removes only a shared, low-rank component of the batch gradient.
The cross-covariance term that carries token-specific semantic information
remains, ensuring that semantic learning is preserved.
This decomposition provides an intuitive bridge between the empirical results in
Section~3 and the formal analysis in Section~4.